\section{Limitations and future work}

The results in \cref{obj:thm:point-frontier,obj:thm:interval-frontier} characterize minimax squared-error risk and honest expected interval length for the finite-cell randomized-MAR class in \cref{obj:def:law-class}, with a known cellwise arrival floor \(q\in[1/2,1]\) and a specified baseline-support bound \(d\). These conditions provide the reference point for several extensions concerning weaker arrival guarantees, adaptation, and practical implementation.

Arrival floors below one half are a natural direction for further analysis. Extending the comparison to \(0<q<1/2\) requires establishing upper and lower bounds that account for the dependence on the floor, including regimes in which it approaches zero as the sample size increases. The bounds established for \(q\in[1/2,1]\) do not determine that dependence. A separate question is adaptation when the floor is unknown. The construction in \cref{obj:def:mm-estimator,obj:def:mm-interval} uses the supplied floor, so a procedure that selects its degree of correction or interval radius from the data would require a new analysis of risk and uniform coverage. Such an analysis would need to address arrival probabilities in sparsely represented cells, where estimating a valid lower bound can itself be difficult.

The comparison with broader models has a precise lower-bound interpretation. Any law class on the same observed experiment, with the same treatment-effect target, that contains the submodels used in \cref{obj:lem:parametric-floor-infrastructure,obj:thm:normalized-converse,obj:lem:prior-reduction} inherits the corresponding lower bounds. For point estimation, the enlarged worst-law risk includes the laws supporting the converse. For honest intervals, coverage over the enlarged class entails coverage over those submodels, and the enlarged worst-law expected length includes their expected lengths. Matching procedures for broader classes require further analysis of their identifying restrictions and statistical complexity; the matching upper bounds here concern the class in \cref{obj:def:law-class}.

Computational implementation is another direction for development. The estimator in \cref{obj:def:mm-estimator} averages over auxiliary randomization conditional on the observed sample. This conditional average specifies a statistical procedure, while an efficient method for evaluating it remains a practical question. Possible approaches include exact evaluation where feasible and numerical approximation elsewhere. An approximation would require control of its additional error, together with an assessment of how that error affects the risk bound and the coverage of intervals centered at the computed estimate. The statistical construction supplies no computational complexity guarantee.

Practical constants also warrant investigation. The polynomial branch of \cref{obj:def:mm-estimator} requires \(n\ge e^{128}-e\), and the proof's fallback contribution is \(C_{\mathrm{fall}}=4e^{128}\). With the exhibited upper-risk constant, the interval in \cref{obj:def:mm-interval} is therefore the full range \([-1,1]\) whenever \(\log(e+n)<128\). The theorems establish universal-constant minimax orders, while these particular numerical choices do not yield an informative interval at ordinary sample sizes. Sharper constants, alternative tuning, and computational study are needed to turn the theoretical construction into a practical procedure; each change would require new risk and coverage justification.

Empirical calibration can complement this work through simulations that vary support size, cell-frequency imbalance, and arrival probabilities within the specified model. Mean squared error, coverage, and expected interval length would provide direct measures of performance, including near complete arrival and in settings with many sparsely occupied cells. Such exercises could guide implementation and identify useful tuning choices, while uniform guarantees for those choices would remain a separate theoretical task.
